\section{Evidence Status}

This appendix preserves the meeting evidence currently archived in the workspace. One early checkpoint is reconstructed from dated project artifacts rather than preserved as a full formal minute set, so it remains explicitly labelled as \textit{recovered}. That distinction is kept visible instead of being hidden.

\begin{table}[H]
\centering
\small
\begin{tabular}{llll}
\toprule
\wcell{3.0cm}{\textbf{Record}} & \wcell{1.9cm}{\textbf{Date}} & \wcell{2.5cm}{\textbf{Type}} & \wcell{4.2cm}{\textbf{Current use}} \\
\midrule
\wcell{3.0cm}{Recovered early implementation checkpoint} & \wcell{1.9cm}{15/03/2026} & \wcell{2.5cm}{Recovered note} & \wcell{4.2cm}{Explains the first playable scope and early implementation decisions.} \\
\wcell{3.0cm}{Kickoff meeting} & \wcell{1.9cm}{17/03/2026} & \wcell{2.5cm}{Formal note} & \wcell{4.2cm}{Shows initial direction-setting and role alignment.} \\
\wcell{3.0cm}{Scope alignment meeting} & \wcell{1.9cm}{18/03/2026} & \wcell{2.5cm}{Formal note} & \wcell{4.2cm}{Shows agreement on tactical-battle scope and documentation format.} \\
\wcell{3.0cm}{Validation planning meeting} & \wcell{1.9cm}{15/04/2026} & \wcell{2.5cm}{Formal note} & \wcell{4.2cm}{Shows testing strategy, defect triage, and final validation planning.} \\
\bottomrule
\end{tabular}
\caption{Meeting evidence preserved in the final-report appendix.}
\label{tab:appendix-meeting-register}
\end{table}

\section{Recovered Early Implementation Checkpoint}

\textbf{Date:} 15/03/2026 \\
\textbf{Type:} recovered from local implementation artifacts \\
\textbf{Attendance:} exact attendance not preserved in the workspace

\subsection{Purpose}

Define the minimum playable slice that could demonstrate the project idea in software rather than only in prose.

\subsection{Recovered Decisions}

\begin{itemize}
    \item Keep one battle scene for the first playable version.
    \item Fix the board at 6x6 so movement, targeting, and occupancy remain readable.
    \item Limit the first playable card set to Move, Strike, Guard, and Push.
    \item Use one shared `UnitController` for player and enemy units.
    \item Keep enemy behaviour deterministic.
    \item Use button-based card selection instead of drag-and-drop.
\end{itemize}

\subsection{Recovered Follow-up}

\begin{itemize}
    \item Implement the first-pass script set and scene bootstrap.
    \item Refresh diagrams so they match the code actually built.
    \item Capture early prototype screenshots.
    \item Start separating validation notes from design notes.
\end{itemize}

\section{Kickoff Meeting}

\textbf{Date:} 17/03/2026 \\
\textbf{Attendance:} Guangjun HU, Hongshuo HU, Zehan CHEN

\subsection{Context}

The team compared initial game ideas, narrowed the project direction, and assigned the first requirement and planning tasks.

\subsection{Decisions}

\begin{itemize}
    \item Confirm \projectname{} as the chosen project direction.
    \item Agree initial role allocation.
    \item Set the first delivery expectations for requirements and design material.
\end{itemize}

\section{Scope Alignment Meeting}

\textbf{Date:} 18/03/2026 \\
\textbf{Attendance:} Guangjun HU, Cheng YE

\subsection{Context}

This meeting aligned the project around a tactical card battle demo and standardised a lightweight documentation format for later work.

\subsection{Decisions}

\begin{itemize}
    \item Keep the software scope focused on a tactical card battle rather than a larger campaign game.
    \item Assign UML and design-diagram production.
    \item Use a repeatable meeting-note structure for later checkpoints.
\end{itemize}

\section{Validation Planning Meeting}

\textbf{Date:} 15/04/2026 \\
\textbf{Format:} Online \\
\textbf{Chair:} Guangjun HU \\
\textbf{Secretary:} Zehan CHEN

\subsection{Attendance}

\begin{itemize}
    \item Guangjun HU -- present
    \item Hongshuo HU -- present
    \item Cheng YE -- present
    \item Zehan CHEN -- present
\end{itemize}

\subsection{Agenda}

\begin{enumerate}
    \item Review usability findings from the playable build.
    \item Confirm automated-test scope and smoke responsibilities.
    \item Allocate drafting tasks for the final validation material.
    \item Set deadlines for open defects and reruns.
\end{enumerate}

\subsection{Key Decisions}

\begin{itemize}
    \item Treat deterministic behaviour as regression-test candidates and keep manual smoke for overlay-heavy UI concerns.
    \item Require rerun-backed evidence before a visible issue can be considered closed.
    \item Keep late work focused on flow, HUD, and readability rather than widening the feature set.
\end{itemize}

\subsection{Action Items}

\begin{table}[H]
\centering
\footnotesize
\setlength{\tabcolsep}{4pt}
\begin{tabular}{llll}
\toprule
\wcell{4.1cm}{\textbf{Task}} & \wcell{2.0cm}{\textbf{Owner}} & \wcell{1.8cm}{\textbf{Due}} & \wcell{2.8cm}{\textbf{Review owner}} \\
\midrule
\wcell{4.1cm}{Draft report sections and supporting tables} & \wcell{2.0cm}{Cheng YE} & \wcell{1.8cm}{16/04/2026} & \wcell{2.8cm}{Guangjun HU} \\
\wcell{4.1cm}{Run EditMode tests and archive results} & \wcell{2.0cm}{Hongshuo HU} & \wcell{1.8cm}{15/04/2026} & \wcell{2.8cm}{Guangjun HU} \\
\wcell{4.1cm}{Update defect register and meeting notes} & \wcell{2.0cm}{Zehan CHEN} & \wcell{1.8cm}{16/04/2026} & \wcell{2.8cm}{Cheng YE} \\
\wcell{4.1cm}{Run final integration check in Unity} & \wcell{2.0cm}{Hongshuo HU} & \wcell{1.8cm}{17/04/2026} & \wcell{2.8cm}{Whole team} \\
\bottomrule
\end{tabular}
\caption{Action items from the validation planning meeting.}
\label{tab:appendix-validation-actions}
\end{table}
